\documentclass{article}
\usepackage{hyperref}
\usepackage{Style}

\nocite{*} % Comentar si quiero citar
%\addbibresource{bibliografia.bib} % Quitar el comentado si quiero usar bibliografia

\begin{document}

\begin{minipage}{2.5cm}
    \includegraphics[width=2cm]{imagen_puc.jpg}
\end{minipage}
\begin{minipage}{12cm}
    {\sc Pontificia Universidad Católica de Chile\\
    Facultad de Matemáticas\\
    Departamento de Matemática\\
    Profesor: Manuel Cabezas -- Ayudante: Benjamín Mateluna}
\end{minipage}
\vspace{2ex}

{\centerline{\bf Introducción a la Geometría - MAT1304}
\centerline{\bf Ayudantía 6 - Repaso I2}}
\centerline{\bf 26 de marzo de 2026}

\vspace{1cm}
\noindent\textbf{Problema 1.} Dado $\triangle ABC$ con $\overline{AB}=6$, $h_{C}=3$ y 
$\overline{BC}=5$, encuentre el largo de $h_{A}$. Donde $h_{X}$ es la altura que pasa por el 
vértice $X$.

\vspace{5mm}
\noindent\textbf{Problema 2.} Encuentre el área de un polígono regular de 6 lados, 
cuyo lado mide a.

\vspace{5mm}
\noindent\textbf{Problema 3.} Dado un rectángulo cualquiera, llamamos $A$ a su área y $p$ a su
perímetro. Demostrar que $p^{2}\geq16A$. ¿Para qué rectángulos se tiene la igualdad?

\vspace{5mm}
\noindent\textbf{Problema 4.} Sea $\triangle ABC$ un triángulo cualquiera. Decimos que 
$a=\overline{BC}$, $b=\overline{AC}$, $c=\overline{AB}$ y $m_{C}$ la mediana trazada desde el 
vértice $C$. Demuestre que
\begin{equation*}
    \frac{a^{2}+b^{2}}{2}=m_{C}^{2}+\frac{c^{2}}{4}
\end{equation*}

\vspace{1cm}
\noindent\textbf{\underline{Ejercicios Propuestos:}}

\vspace{5mm}
\noindent\textbf{Extra 1.} Sean $P_{1}$, $P_{2}$, $P_{3}$, $P_{4}$, $P_{5}$ y $P_{6}$, los 
vértices de un polígono convexo de lados $\overline{P_{1}P_{2}}$, $\overline{P_{2}P_{3}}$, 
$\overline{P_{3}P_{4}}$, $\overline{P_{4}P_{5}}$, $\overline{P_{5}P_{6}}$ y 
$\overline{P_{6}P_{1}}$. Asuma que  cada lado mide 1 y que los ángulos interiores del polígono son
iguales.
\begin{enumerate}
    \item[a)] Encontrar el área de $\triangle P_{1}P_{4}P_{5}$.
    \item[b)] Encontrar el área de $\triangle P_{1}P_{2}P_{3}$. 
\end{enumerate}

\vspace{5mm}
\noindent\textbf{Extra 2.} Demuestre que las medianas de un triángulo cualquiera son 
concurrentes.

%\printbibliography % Quitar el comentado si quiero usar bibliografia

\end{document}
